% ACL requires an unnumbered "Limitations" section (excluded from the page limit).
\section*{Limitations}
\label{sec:limitations}

\paragraph{Single student-model family and scale.} All trained student policies
use Qwen3.5-4B. Generalization across student-model families and to larger scales
(e.g.\ 8B) is future work.

\paragraph{Benchmarks.} SGCD is evaluated on two benchmarks (AppWorld and
$\tau^3$-airline). Broader coverage across additional agentic domains is future
work.

\paragraph{External-model dependence.} The credit reference is
produced by \creditreferenceproviderphrase{} during training. Its cost,
availability, and quality are practical constraints. The inference-time student has no
such dependency, but training-time quality still matters because poor credit
references can produce weak or noisy credit features. We do not claim that the
reference model choice is immaterial. The credit-reference model family and the
structured prompt used to elicit it are described in \suppref{sec:appendix}. The
$\tau^3$ data generation, training environment, and held-out evaluation also use
an external instruction-following user simulator shared across compared arms;
this is an evaluation dependency, not a component of the deployed student.

\paragraph{Credit reference is not formal ground truth.} SGCD approximates
stepwise reference behavior from verified siblings. This is weaker than a
benchmark-provided stepwise oracle and can miss valid alternative plans. The
bounded credit weights reduce the risk of overcommitting to one interpretation,
but they do not make the credit reference infallible.

\paragraph{Training-run coverage.} AppWorld values are single-run point estimates because time and compute constraints prevented multi-seed reruns, so we make no AppWorld variance or significance claim. On $\tau^3$-airline, SGCD uses three independent training runs (Table~\ref{tab:seedrobust}), but Vanilla GRPO is a fixed single-run anchor; the reported dispersion is therefore SGCD-only and the comparison remains directional rather than variance-controlled.

\paragraph{AppWorld recipe comparison.} The AppWorld arms match rollout,
environment, and evaluation settings, but not their regularization term:
only Vanilla GRPO includes reference-KL. SGCD instead uses detached dense
divergence for credit weighting, so the result compares complete recipes and
does not isolate the effect of removing reference-KL.

\paragraph{Diagnostic and ablation scope.} The AppWorld SFT regression table is
a diagnostic subset result, not a headline held-out result. The full SGCD result
also depends on the exact external-LLM prompt and dynamic-sampling rollout set;
we report the core controls but do not exhaustively sweep every prompt-builder
or divergence variant.

\paragraph{Artifact availability.} The review package does not include the full
training/evaluation code, generated training trajectories and corrected task
artifacts, or exact run-level manifests. The supplement reports the available
settings, prompt template, and evaluation definitions, but the submitted
materials do not support exact end-to-end replication.
